\section{Introduction}
\label{sec:intro}

Artificial intelligence has assisted scientific research for decades. Until recently, "AI for science" largely meant a specialized model targeting one well-defined task—predicting biomolecular structures\citep{jumper2021highly,krishna2024generalized,lin2023evolutionary,chai2024chai,wohlwend2025boltz,passaro2025boltz,wu2022high,evans2021protein}, accelerating materials discovery\cite{merchant2023scaling,zeni2025generative,xie2018crystal,deng2023chgnet,batatia2022mace}, or forecasting the weather\cite{lam2023learning,bi2022pangu}—while human researchers still framed the question, built the AI tools themselves, interpreted the results, and communicated the findings. AlphaFold\cite{jumper2021highly} epitomizes this paradigm: a landmark solution to a single, long-standing problem, but one confined to that problem alone. The rapid advance of large language models and agentic scaffolds is now reshaping this picture. Combining foundation models, external tools, and agentic workflows, a single system can carry a study from an initial hypothesis through literature review, experimentation, and analysis to a written draft—a paradigm now referred to as \emph{AutoResearch}\cite{tie2026autoresearch,kong2026ai,zhang2026far}. AI is gradually becoming a participant in scientific discovery rather than merely a tool for it. What is being automated is no longer an isolated step, but the research process as a whole. A natural question follows: \textbf{how well can agents actually perform across the full AutoResearch process?}



Evaluating autonomous research agents requires characterizing their behavior across the entire scientific process. Recent work has made notable strides toward this goal, from benchmarks that assess isolated research skills \citep{rein2023gpqa, phan2025humanity, tian2024scicode, siegel2024core, hu2025repro, starace2025paperbench} to increasingly ambitious end-to-end pipelines \citep{wang2026naturebench, xu2026researchclawbench, bragg2026astabench}. Despite such progress, existing evaluations leave three gaps that together obscure \emph{how} an agent works and \emph{where} it breaks down.



\textbf{Tasks are narrowly-scoped.} Existing discovery benchmarks fall into
two groups, and neither captures the full research lifecycle in real scientific
domains. The first places agents in simulated or fictional worlds, adopted
because real experiments are expensive---but this comes at the cost of
realism~\citep{jansen2024discoveryworldvirtualenvironmentdeveloping,
wang2022scienceworldagentsmarter5th, majumder2024discoverybenchdatadrivendiscoverylarge,
li2026pdagentbenchcharacterizinggroundingarchitecting}. The second stays
grounded in real science but restricts itself to verifiable settings---most
commonly machine learning or software engineering tasks whose outcomes can be
measured against a well-defined target~\citep{chan2025mle, nathani2025mlgym,
starace2025paperbench, chen2025scienceagentbenchrigorousassessmentlanguage,
tian2024scicode, wang2026naturebench, xu2026researchclawbench}. Coverage of
science domains is correspondingly narrow, and genuine open-ended discovery
tasks are excluded.

\textbf{Evaluation measures performance, not process.} Nearly all existing
benchmarks anchor scoring at the endpoint---a reference match, a reproduced
result, or a published SOTA. This invites reward hacking: circular validation,
grader-fitting, or leakage move the number without doing the science, and the
score cannot distinguish a sound trajectory from a gamed
one~\citep{chan2025mle, hu2025repro, wang2026naturebench}. It also compresses a
long-horizon trajectory into a single scalar, reporting \emph{that} a run
failed but not \emph{why}, \emph{where}, or \emph{how}---so endpoint
evaluation can rank systems but cannot diagnose them.

\textbf{Failure diagnoses lack systematic coverage or artifact-level
visibility.} Expert case studies scrutinize individual trajectories in
depth~\citep{kirgis2026shadow, li2026scicomp, rawat2026plausible,
horstmann2026neuro, eulig2026cawm, trehan2026llmsarentscientistsyet,
zhang2026far} but cannot generalize across systems or tasks.
Categorization efforts targeting scientific
discovery~\citep{luo2025automateseehiddenpitfalls,
eulig2026cawm, bisht2026agenticaiscientistsbuilt} offer
broader categories but are derived from small corpora. The most systematic
line---trace-level analysis of multi-agent or QA
settings~\citep{cemri2026multi, deshpande2025trailtracereasoningagentic, zhang2025agent}---annotates
conversational traces, so artifact-level failures stay invisible and stages are
interaction phases rather than steps of the scientific method.

Together these gaps leave a critical blind spot: \textbf{current evaluations may show that an agent succeeded, but reveal little about how it operates, or precisely where it breaks down.}




To address this gap, we investigate why autonomous research agents fail, producing two artifacts. The first is \textbf{\bench}, a collection of $800$ research trajectories with \textbf{artifact-aware, process-level failure annotations}, elicited from eight harness--model combinations on a $100$-task suite spanning the \textbf{full-lifecycle} scientific workflow---ideation, retrieval \& synthesis, execution, analysis, writing, and review---across \textbf{seven frontier domains} and constructed from papers in prestigious venues. Tasks span two regimes: those with an explicit execution-feedback signal (a human SOTA or quantitative metric), and fully open-ended tasks with none, where we deliberately judge the rigor and self-consistency of the agent's \emph{process} rather than agreement with ground truth. Annotations come from a \textbf{human-calibrated, artifact-aware agent-as-a-judge annotator} \citep{zhuge2024agentjudge} that reads each trajectory in full---run logs, generated data, code, and reports---rather than its final answer. We calibrate it against human annotations and spot-check its outputs post hoc.

To enable systematic annotation and comprehensive analysis of \bench, we develop the
second artifact: \textbf{\taxonomy}, the first systematic and comprehensive
failure taxonomy for autonomous research agents. \taxonomy is empirically
grounded in the annotated trajectories and organizes $45$ failure patterns
on two cross-cutting axes: the lifecycle \emph{stage} at which a failure
manifests, and the underlying \emph{root cause}, so that superficially
different failures sharing a mechanism are grouped together
(Figure~\ref{fig:failure-attribution}). Together, \bench supplies the
empirical evidence of how research agents fail in practice, while \taxonomy
provides the structured vocabulary to diagnose, attribute, and ultimately
mitigate these failures.

Our analysis of \bench yields one core finding that revises prevailing
assumptions. While failure patterns span all stages of the research
lifecycle, they converge on a single overarching limitation: current agents
lack a \textbf{metacognitive loop}---the ability to check what they produced
against what they found, revise when it does not hold up, and question whether
the path they took was sound.


\paragraph{Contributions.}
In summary, our core contributions are as follows:

\begin{itemize}
    \item \textbf{\bench}, a \textbf{open-source} collection of $800$ trajectories with
    artifact-aware, process-level failure annotations, released together
    with the underlying $100$-task, seven-domain, full-lifecycle task suite
    built from published frontier science, including a fully open-ended
    subset scored on process rather than outcome;
    
    \item \textbf{\taxonomy}, the first systematic failure taxonomy for autonomous research agents, organizing 45 empirically-grounded failure patterns;

    \item \textbf{A human-calibrated, artifact-aware agent-as-a-judge}
    that is used for \taxonomy categorization. It reads each trajectory in full rather than its final answer alone and hence supporting analyzing and understanding failure patterns;

    \item \textbf{A systematic diagnosis of when and why these agents fail},
    tracing all three cognitive root causes to a shared \textbf{metacognitive
    deficit}---the absence of a closed \textbf{metacognitive loop}. 
    This diagnosis locates the gap at the model level---the same patterns recur
across all eight harness--model combinations---and identifies the metacognitive
loop as a capacity that next-generation language models will need for genuine
scientific discovery. Whether orchestration alone can compensate is an open
question we do not test.

    

  
\end{itemize}